\documentclass[a4paper,12pt]{article}

\usepackage[a4paper,margin=1in]{geometry}
\usepackage{hyperref}

\title{Computer Science Portfolio Project Report}
\author{
Sayed Mohammad Kamran Abbas Haider\\
Student Number: GH1046612
}
\date{\today}

\begin{document}

\maketitle

\tableofcontents

\newpage

\section{Introduction}

This report explains the development of my Computer Science Portfolio website. I created
this portfolio as part of the Computer Science Lab module at GISMA University of Applied
Sciences. The main purpose of the website is to show my technical skills, projects and
personal information in one place.

The portfolio contains an About Me section, a Skills section, a Projects section and a
Contact section. It also includes a downloadable Curriculum Vitae (CV) that I created
using LaTeX. During this project I improved my understanding of HTML, CSS, Git,
GitHub, GitHub Pages and LaTeX.

This report explains how my portfolio is organised, the technologies that I used,
the design decisions I made while creating the website and what I learned during the
project. I also discuss the strengths, weaknesses and possible improvements for the
future.

\section{Portfolio Website and GitHub URLs}

GitHub Profile

\url{https://github.com/1granny1}

\vspace{0.5cm}

Portfolio Website

\url{https://1granny1.github.io/Portfolio-For-Computer-Science-lab/}

\vspace{0.5cm}

Portfolio GitHub Repository

\url{https://github.com/1granny1/Portfolio-For-Computer-Science-lab}

\vspace{0.5cm}

University Database System Repository

\url{https://github.com/1granny1/University-Database-System}

\vspace{0.5cm}

Snake Game Repository

\url{https://github.com/1granny1/Snake-game-for-python}

```latex
\section{Portfolio Structure}

\subsection{Website Structure}

My portfolio website is organised into different sections to make it easy to
navigate. At the top of the page there is a navigation bar with links to the
Home, About, Skills, Projects and Contact sections.

The Home section introduces me as a Software Engineering student at GISMA
University of Applied Sciences. It also contains my profile picture and a
button to download my CV.

The About section gives a short introduction about myself and explains my
interest in programming and software development.

The Skills section lists the programming languages and tools that I have
learned so far, including HTML, CSS, Python, Java, SQL, Git, GitHub and
LaTeX.

The Projects section contains three projects. The first project is my
Computer Science Portfolio website. The second project is my University
Database System developed using SQL. The third project is a Snake Game
developed using Python. Each project contains a short description and a
link to its GitHub repository.

The Contact section contains my email address and a link to my GitHub
profile so that visitors can easily contact me or view my work.

\subsection{GitHub Repository Structure}

My GitHub repository contains all the files needed to run my portfolio
website. The main file is \texttt{index.html}, which contains the content
of the website. The design of the website is stored in
\texttt{style.css}.

The repository also contains my profile picture, my CV in PDF format and
the README file. GitHub Pages is used to host the website online so that
it can be viewed using a web browser.

\section{Design Decisions}

When creating my portfolio, I wanted to keep the design simple and easy to
use. I chose a blue and white colour theme because it looks clean and
professional.

I organised the website into different sections so visitors can quickly
find information without opening different pages. I also added a
navigation bar to make moving around the website easier.

The projects section was designed to clearly show the work I have
completed during this module. I also included a button to download my CV
so employers or lecturers can easily view it.

\section{User Experience}

I wanted the website to be simple for anyone to use. The navigation links
allow visitors to move to different sections of the page quickly.

The text is easy to read because I used clear headings, simple colours and
enough spacing between sections. The project links take users directly to
my GitHub repositories, where they can view my work.

\section{Technological Choices}

I used HTML to create the structure of the website and CSS to improve its
appearance. Git and GitHub were used to keep track of changes and manage
my files.

GitHub Pages was used to publish my portfolio online because it is free
and easy to use. I created my CV using LaTeX because it produces a clean
and professional document.

\section{Tools and Technologies}

\begin{itemize}

\item HTML

\item CSS

\item Python

\item SQL

\item Git

\item GitHub

\item GitHub Pages

\item LaTeX

\item Visual Studio Code

\end{itemize}

```latex
\section{Reflection and Future Work}

\subsection{Strengths}

I think one of the main strengths of my portfolio is that it is simple and
easy to use. Visitors can quickly find information about me, my skills and
my projects without any confusion. The website is organised into clear
sections, which makes navigation easy.

Another strength is that I uploaded the website using GitHub Pages, so it
can be accessed from anywhere using a web browser. I also included my CV
and links to my GitHub repositories, making it easier for people to see my
work.

\subsection{Weaknesses}

Although I am happy with the final result, there are still some things that
could be improved. The website is quite simple and does not have many
interactive features. I also think some sections could contain more
information about my projects.

Another weakness is that I am still learning web development, so there are
many design techniques that I have not used yet.

\subsection{Future Improvements}

In the future I would like to improve my portfolio by adding more projects
as I continue my Software Engineering course.

I also want to improve the design by learning more CSS and making the
website look more modern. Another improvement would be adding more
information about each project, including screenshots and a longer
description.

As I learn JavaScript and backend development, I would like to make my
portfolio more interactive and add new features.

\section{Conclusion}

This portfolio project helped me improve my understanding of HTML, CSS,
Git, GitHub, GitHub Pages and LaTeX. During this project I learned how to
create a personal portfolio website, upload it to GitHub and organise my
projects in one place.

I believe this portfolio is a good starting point for my future work. I
plan to continue updating it by adding new projects and improving the
design as I learn more programming skills throughout my degree.

\end{document}
```
